\begin{corollary}[7.6.2]
\label{IV.2.s7.6.2-fr}
Bajo las hipótesis de \textup{(7.6.1)}, existe una correspondencia biyectiva
canónica entre el conjunto de los ideales maximales $\mathfrak m_j$
de $A'$ \textup{(es decir, el conjunto de los puntos de
$\operatorname{Spec}(A')$ situados sobre el punto cerrado de $A$)} y
el conjunto de los ideales primos minimales $\mathfrak q_j$ de
$\widehat A$ \textup{(es decir, el conjunto de los puntos maximales de
$\operatorname{Spec}(\widehat A)$)}; en esta correspondencia, el
completado $\widehat{A}'_{\mathfrak m_j}$ de $A'_{\mathfrak m_j}$
es isomorfo al cierre entero de $\widehat A/\mathfrak q_j$.
\end{corollary}

\begin{proof}
En efecto, se sabe que el cierre entero de $\widehat A$ en su anillo
total de fracciones es el producto directo de los cierres enteros de
los $\widehat A/\mathfrak q_j$, que son anillos locales completos
(\textbf 0, 23.1.6).
\end{proof}
